\documentclass{article}

\usepackage{amsmath,amssymb}
\usepackage{xurl}
\usepackage[hidelinks]{hyperref}

\input{command}

\title{Wolf Theorem~6.14: Fixed-Point Classification Completion Record}
\author{}
\date{2026-07-19}

\begin{document}
\maketitle
\gapnote{open-gap}{resolved}

\section{The retraction classification}

Proposition~1.5 of Ref.~\cite{Wolf2012Quantum} classifies a positive linear
retraction onto a unital finite-dimensional $*$-subalgebra; its block formula
appears in source Equation~(1.40). The local source is
\path{Notes/WolfNoteTexSource/ch01_deconstructing_quantum.tex},
lines~536--570.

The formal theorem
\leanid{StarSubalgebra.exists_block_densities_of_positive_retraction} proves
the complete proposition. It assumes that the map is positive, that its image
lies in the subalgebra, and that it fixes every element of the subalgebra. The
theorem constructs the unitary direct-sum representation, uses positivity to
eliminate off-diagonal inputs, and produces a positive semidefinite trace-one
matrix for each summand. It adds no complete-positivity, trace-preservation, or
Schwarz hypothesis.

\section{The fixed-point theorem}

Let $T:\MN{D}\to\MN{D}$ be positive and trace preserving, and let $T^*$ denote
its trace-pairing adjoint. Assume that $T^*$ satisfies the Schwarz inequality.
Theorem~6.14 of Ref.~\cite{Wolf2012Quantum} asserts the existence of a unitary
$U$, a decomposition
\begin{align}
  \mathbb C^D
    &=\mathcal H_0
      \oplus\bigoplus_{k=1}^K
      \mathbb C^{d_k}\otimes\mathbb C^{m_k},
  \label{eq:wolf_theorem6_14_fixed_point_projection_gap_source_decomposition}
\end{align}
and positive-definite density matrices
$\rho_k\in M_{m_k}(\mathbb C)$ such that source Equation~(6.63) is
\begin{align}
  \operatorname{Fix}(T)
    &=U\left(0\oplus\bigoplus_{k=1}^K
      M_{d_k}(\mathbb C)\otimes\rho_k\right)U^*.
  \label{eq:wolf_theorem6_14_fixed_point_projection_gap_source_fixed_space}
\end{align}
The summand $\mathcal H_0$ is a genuine zero summand: every fixed point
vanishes on it. The local source for Theorem~6.14 is
\path{Notes/WolfNotePDF/wolf_ch6_sections.txt}, lines~896--905.

\section{Mean-ergodic retraction}

Every positive trace-preserving matrix endomorphism has bounded forward
orbits. Its mean-ergodic projection is positive and trace preserving, is
idempotent, and has range exactly $\operatorname{Fix}(T)$. If $T$ is unital,
the projection is also unital.\footnote{Formal declarations:
\leanid{IsPositiveMap.hasBoundedOrbits_of_tracePreserving},
\leanid{IsPositiveMap.meanErgodicProjection_isPositiveMap}, and
\leanid{IsTracePreservingMap.meanErgodicProjection_isTracePreservingMap}.}

The trace-pairing adjoint of the mean-ergodic projection is a positive unital
idempotent retraction onto $\operatorname{Fix}(T^*)$.\footnote{Formal
declarations:
\leanid{LinearMap.HasBoundedOrbits.range_traceAdjointMap_meanErgodicProjection}
and
\leanid{IsPositiveMap.traceAdjoint_meanErgodicProjection_isPositiveUnitalRetraction}.}
The proof uses the complementary decomposition for the original mean-ergodic
projection and nondegeneracy of the trace pairing. It requires no
bounded-orbit or spectral hypothesis for $T^*$.

\section{Full-support block retraction}

Assume first that $T$ has a positive-definite fixed point. The Schwarz equality
argument makes $\operatorname{Fix}(T^*)$ a unital $*$-algebra, and the adjoint
mean-ergodic projection is a positive retraction onto this algebra. Applying
Proposition~1.5 of Ref.~\cite{Wolf2012Quantum} gives the weighted
partial-trace block form of the projection.\footnote{Formal declaration:
\leanid{IsPositiveMap.exists_block_densities_of_adjoint_meanErgodicProjection}
in
\doclink{QICLean/Channel/FixedPoint/FullSupportBlockRetraction}.}
This formalizes the conditional-expectation step at local source
lines~1483--1492.

The formalization orders each block as
$\mathbb C^{m_k}\otimes\mathbb C^{d_k}$, so the adjoint fixed-point algebra is
$\Id_{m_k}\otimes M_{d_k}(\mathbb C)$. The partial trace is therefore taken
over the first, multiplicity factor:
\begin{align}
  P_T(B)
    &=U\left(\bigoplus_k
      \sigma_k\otimes
      \operatorname{tr}_{m_k}((U^*BU)_{kk})\right)U^*.
  \label{eq:wolf_theorem6_14_fixed_point_projection_gap_projection_block_form}
\end{align}
Using $\operatorname{tr}_{d_k}$ in
(\ref{eq:wolf_theorem6_14_fixed_point_projection_gap_projection_block_form})
would reverse the factor convention and would be dimensionally inconsistent.
The range is
\begin{align}
  \operatorname{ran}(P_T)
    &=U\left(\bigoplus_k
      \sigma_k\otimes M_{d_k}(\mathbb C)\right)U^*,
  \label{eq:wolf_theorem6_14_fixed_point_projection_gap_full_support_range}
\end{align}
and therefore equals $\operatorname{Fix}(T)$ under the full-support
hypothesis.\footnote{Formal declaration:
\leanid{IsPositiveMap.exists_block_densities_of_meanErgodicProjection}
in
\doclink{QICLean/Channel/FixedPoint/TraceAdjointDensityBlocks}.}

\section{Maximal stationary support}

For every positive trace-preserving map, define
\begin{align}
  \rho_0
    &=T_\infty(\Id).
  \label{eq:wolf_theorem6_14_fixed_point_projection_gap_maximal_support_state}
\end{align}
The matrix $\rho_0$ is positive semidefinite, and its support contains both the
support and the range of every fixed point.\footnote{Formal declaration:
\leanid{IsPositiveMap.exists_maximalSupport_fixedPoint} in
\doclink{QICLean/Channel/FixedPoint/MaximalSupportBasic}.}
This supplies the maximal-support witness in the full generality of the source.

Let $\rho\succeq0$ be any fixed point and let $Q$ be its support projection.
Every positive matrix in the corner $Q\MN{D}Q$ is mapped back into that corner.
Decomposing an arbitrary matrix into four positive matrices extends the result
to the entire corner. Choose an isometry $V$ satisfying
\begin{align}
  V^*V
    &=\Id,
  \label{eq:wolf_theorem6_14_fixed_point_projection_gap_isometry_initial}\\
  VV^*
    &=Q.
  \label{eq:wolf_theorem6_14_fixed_point_projection_gap_isometry_range}
\end{align}
The compressed map is
\begin{align}
  \widetilde T(X)
    &=V^*T(VXV^*)V.
  \label{eq:wolf_theorem6_14_fixed_point_projection_gap_compressed_map}
\end{align}
It is positive and trace preserving and has a positive-definite fixed point.
It also satisfies source Equation~(6.52) of Ref.~\cite{Wolf2012Quantum}:
\begin{align}
  T(VXV^*)
    &=V\widetilde T(X)V^*.
  \label{eq:wolf_theorem6_14_fixed_point_projection_gap_compression_intertwining}
\end{align}

When $\rho=\rho_0$, maximality of the support gives source Equation~(6.51) of Ref.~\cite{Wolf2012Quantum}:
\begin{align}
  \operatorname{Fix}(T)
    &=V\operatorname{Fix}(\widetilde T)V^*.
  \label{eq:wolf_theorem6_14_fixed_point_projection_gap_compressed_fixed_space}
\end{align}
Thus every ambient fixed point vanishes on the complementary block.

For a general stationary $\rho$, the auxiliary fixed-space equivalence assumes
that $QXQ=X$ for every $X\in\operatorname{Fix}(T)$. This is the hypothesis
named \texttt{hmax} in
\leanid{IsPositiveMap.fixedPoint_iff_exists_fixedPoint_stationarySupportCompression}.
The source-faithful theorem
\leanid{IsPositiveMap.exists_maximalSupportCompression} exposes no such
hypothesis: it sets $\rho_0=T_\infty(\Id)$ and derives support maximality from
\leanid{IsPositiveMap.exists_maximalSupport_fixedPoint}.\footnote{Formal
declarations:
\leanid{IsPositiveMap.map_supported_on_fixedPoint_support} and
\leanid{IsPositiveMap.exists_maximalSupportCompression}
in
\doclink{QICLean/Channel/FixedPoint/StationarySupportRestriction}.}

\section{Schwarz compression and positive-definite blocks}

The trace adjoint commutes with isometric compression:
\begin{align}
  \widetilde T^*(A)
    &=V^*T^*(VAV^*)V.
  \label{eq:wolf_theorem6_14_fixed_point_projection_gap_compressed_trace_adjoint}
\end{align}
The Schwarz defect of $\widetilde T^*$ is the compression of the Schwarz defect
of $T^*$ plus a positive term supported on $\Id-VV^*$. Therefore, the Schwarz
hypothesis passes from $T^*$ to $\widetilde T^*$ using only the positivity of
$T$ already assumed in Theorem~6.14 of Ref.~\cite{Wolf2012Quantum}.

The full-support theorem applied to $\widetilde T$ gives
\begin{align}
  \operatorname{Fix}(\widetilde T)
    &=U\left(\bigoplus_k
      \sigma_k\otimes M_{d_k}(\mathbb C)\right)U^*,
  \label{eq:wolf_theorem6_14_fixed_point_projection_gap_compressed_density_blocks}
\end{align}
where each $\sigma_k$ is positive semidefinite and has trace one. Let
$\rho_{\mathrm{full}}$ be the positive-definite fixed point retained by the
compression theorem. Then
\begin{align}
  U^*\rho_{\mathrm{full}}U
    &=\bigoplus_k\sigma_k\otimes X_k.
  \label{eq:wolf_theorem6_14_fixed_point_projection_gap_full_support_state_blocks}
\end{align}
Every principal block $\sigma_k\otimes X_k$ is positive definite. Taking the
two partial traces gives
\begin{align}
  \operatorname{tr}_{m_k}(\sigma_k\otimes X_k)
    &=X_k,
  \label{eq:wolf_theorem6_14_fixed_point_projection_gap_partial_trace_multiplicity}\\
  \operatorname{tr}_{d_k}(\sigma_k\otimes X_k)
    &=\tr(X_k)\sigma_k.
  \label{eq:wolf_theorem6_14_fixed_point_projection_gap_partial_trace_matrix}
\end{align}
The first identity implies that $X_k$ is positive definite, so
$\tr(X_k)>0$. The second then implies that $\sigma_k$ is positive definite.
This argument changes no block dimension and introduces no additional kernel
summand.\footnote{Formal declarations:
\leanid{Matrix.traceAdjointMap_stationarySupportCompression},
\leanid{IsSchwarzMap.traceAdjointMap_stationarySupportCompression}, and
\leanid{IsPositiveMap.exists_block_densities_of_maximalSupportCompression}
in
\doclink{QICLean/Channel/FixedPoint/SupportCompressedDensityBlocks}. The
positive-definite partial-trace results are
\leanid{Matrix.PosDef.partialTraceRight},
\leanid{Matrix.PosDef.partialTraceLeft}, and
\leanid{Matrix.PosDef.left_of_kronecker_of_trace_eq_one} in
\doclink{QICLean/Channel/PartialTrace}.}

\section{Ambient transport}

Let $U_c$ be the compressed block unitary and define
\begin{align}
  W
    &=VU_c.
  \label{eq:wolf_theorem6_14_fixed_point_projection_gap_compressed_isometry}
\end{align}
The map $W$ is an isometry from the compressed space into the original space.
Compare it with the canonical inclusion of $\mathbb C^n$ into
$\mathbb C^{D-n}\oplus\mathbb C^n$. The two inclusions have the same Gram
matrix. A finite-dimensional unitary extension therefore gives an ambient
unitary $U$ satisfying
\begin{align}
  WAW^*
    &=U(0_{D-n}\oplus A)U^*
    \qquad\text{for every }A\in M_n(\mathbb C).
  \label{eq:wolf_theorem6_14_fixed_point_projection_gap_unitary_zero_extension}
\end{align}
Combining
(\ref{eq:wolf_theorem6_14_fixed_point_projection_gap_compressed_fixed_space}),
(\ref{eq:wolf_theorem6_14_fixed_point_projection_gap_compressed_density_blocks}),
and
(\ref{eq:wolf_theorem6_14_fixed_point_projection_gap_unitary_zero_extension})
gives
\begin{align}
  \operatorname{Fix}(T)
    &=U\left(0_{D-n}\oplus
      \bigoplus_k\sigma_k\otimes M_{d_k}(\mathbb C)\right)U^*.
  \label{eq:wolf_theorem6_14_fixed_point_projection_gap_ambient_fixed_space}
\end{align}
The first block is exactly the orthogonal complement of maximal stationary
support. It is the only zero summand because every $\sigma_k$ is positive
definite.\footnote{Formal declarations:
\leanid{Matrix.exists_unitary_zero_extension_eq} in
\doclink{QICLean/Algebra/MatrixGramUnitary}, and
\leanid{IsPositiveMap.exists_fixedPoints_densityBlocks_with_zero} in
\doclink{QICLean/Channel/FixedPoint/WolfTheorem614}.}

\section{Conclusion}

Wolf's Theorem~6.14 is now formalized with its full source hypotheses and
conclusion \cite[Theorem~6.14]{Wolf2012Quantum}. The theorem assumes exactly
that $T$ is positive and trace preserving and that $T^*$ satisfies the Schwarz
inequality. Its fixed-point space has the single complementary zero summand
asserted in source Equation~(6.63). This document therefore records a resolved
paper gap.

\bibliographystyle{alpha}
\bibliography{references}

\end{document}
